% ============================================================
%  DINÁMICA DE SISTEMAS — ÚLTIMA MILLA Y CONGESTIÓN URBANA
%  Santiago de Chile
%  Compilar: pdflatex -shell-escape ultima_milla_santiago.tex
% ============================================================
\documentclass[12pt,a4paper]{article}

% ── Encoding & Language ──────────────────────────────────────
\usepackage[utf8]{inputenc}
\usepackage[T1]{fontenc}
\usepackage[spanish,es-nodecimaldot]{babel}

% ── Layout ───────────────────────────────────────────────────
\usepackage[top=2.5cm,bottom=2.5cm,left=2.8cm,right=2.8cm]{geometry}
\usepackage{setspace}
\setstretch{1.15}
\usepackage{parskip}

% ── Math ─────────────────────────────────────────────────────
\usepackage{amsmath,amssymb,mathtools}

% ── Graphics & Diagrams ──────────────────────────────────────
\usepackage{graphicx}
\usepackage{tikz}
\usetikzlibrary{arrows.meta,positioning,shapes.geometric,
                decorations.pathmorphing,calc,fit,backgrounds,
                arrows,shapes.arrows,bending}
\usepackage{pgfplots}
\pgfplotsset{compat=1.18}

% ── Tables ───────────────────────────────────────────────────
\usepackage{booktabs,longtable,array,tabularx,multirow}
\usepackage{colortbl}

% ── Colors ───────────────────────────────────────────────────
\usepackage[dvipsnames]{xcolor}
\definecolor{azulUCH}{RGB}{0,68,124}
\definecolor{verdeSD}{RGB}{0,128,96}
\definecolor{rojoSD}{RGB}{180,40,40}
\definecolor{ambarSD}{RGB}{200,140,0}
\definecolor{grisClaro}{RGB}{245,245,245}
\definecolor{grisOscuro}{RGB}{80,80,80}

% ── Hyperlinks ───────────────────────────────────────────────
\usepackage[colorlinks=true,
            linkcolor=azulUCH,
            citecolor=verdeSD,
            urlcolor=azulUCH]{hyperref}
\usepackage{url}

% ── Listings (code) ──────────────────────────────────────────
\usepackage{listings}
\lstset{basicstyle=\ttfamily\small,breaklines=true,
        frame=single,backgroundcolor=\color{grisClaro}}

% ── Captions & Floats ────────────────────────────────────────
\usepackage{caption,subcaption,float}
\captionsetup{labelfont=bf,font=small}

% ── Headers / Footers ────────────────────────────────────────
\usepackage{fancyhdr}
\pagestyle{fancy}
\fancyhf{}
\fancyhead[L]{\small\textcolor{grisOscuro}{Dinámica de Sistemas}}
\fancyhead[R]{\small\textcolor{grisOscuro}{Última Milla — Santiago}}
\fancyfoot[C]{\thepage}
\renewcommand{\headrulewidth}{0.4pt}

% ── Boxes ────────────────────────────────────────────────────
\usepackage{tcolorbox}
\tcbuselibrary{skins,breakable}
\newtcolorbox{definicion}[1][]{
  colback=azulUCH!5,colframe=azulUCH!60,
  fonttitle=\bfseries,title=Definición,#1}
\newtcolorbox{arquetipobox}[2][]{
  colback=verdeSD!5,colframe=verdeSD!60,
  fonttitle=\bfseries,title=Arquetipo: #2,#1}
\newtcolorbox{databox}[1][]{
  colback=ambarSD!5,colframe=ambarSD!60,
  fonttitle=\bfseries,title=Fuentes de datos,#1}

% ── Bibliography ─────────────────────────────────────────────


% ── Misc ─────────────────────────────────────────────────────
\usepackage{enumitem}


% ============================================================
\begin{document}

% ── Portada ──────────────────────────────────────────────────
\begin{titlepage}
\centering
\vspace*{1cm}
{\color{azulUCH}\rule{\linewidth}{1.5pt}}\\[0.8em]
{\LARGE\bfseries Dinámica de Sistemas:\\[0.3em]
Última Milla y Congestión Urbana\\[0.3em]
en Santiago de Chile}\\[0.6em]
{\color{azulUCH}\rule{\linewidth}{1.5pt}}

\vspace{2cm}
{\large Modelamiento Causal, Arquetipos y\\
Ecuaciones Diferenciales del Sistema Logístico Urbano}

\vspace{3cm}
\begin{tabular}{ll}
\textbf{Curso:}   & Modelamiento de Sistemas \\[4pt]
\textbf{Semestre:} & 1\textsuperscript{er} semestre \\[4pt]
\textbf{Área:}    & Logística Urbana / Ingeniería de Sistemas \\[4pt]
\textbf{Región:}  & Metropolitana de Santiago, Chile \\[4pt]
\textbf{Fecha:}   & \today \\
\end{tabular}

\vfill
{\small\textcolor{grisOscuro}{Documento generado para uso académico.
Los datos corresponden a fuentes oficiales chilenas.}}
\end{titlepage}

\tableofcontents
\newpage

% ============================================================
\section{Introducción y Motivación}
% ============================================================

El crecimiento sostenido del comercio electrónico en Chile — con un
incremento de la facturación del 30\,\% anual entre 2018 y 2023 según
la Cámara de Comercio de Santiago (CCS) — ha multiplicado la cantidad
de vehículos de reparto circulando en la Región Metropolitana.
Este fenómeno genera una dinámica de retroalimentación no lineal:
más paquetes exigen más vehículos, más vehículos generan más
congestión, la congestión aumenta los tiempos y costos de entrega,
y esos costos presionan a los operadores a añadir más vehículos
para compensar — cerrando un bucle de refuerzo que se autoalimenta.

Al mismo tiempo, operan mecanismos de balance: el alza en el costo
por paquete eleva el precio al consumidor, reduciendo la demanda de
ecommerce en segmentos de menor ingreso; la regulación municipal
sobre zonas de carga y descarga limita el acceso de vehículos;
y la adopción de modelos alternativos (bicicletas de carga,
casilleros inteligentes, consolidación de rutas) puede absorber
parte del flujo.

Este documento desarrolla un modelo de dinámica de sistemas completo
para este problema, incluyendo:
\begin{itemize}[leftmargin=1.5cm]
  \item Descripción del sistema y variables relevantes
  \item Diagrama causal (bucles de retroalimentación)
  \item Diagrama de stocks y flujos
  \item Arquetipos sistémicos identificados
  \item Ecuaciones diferenciales del modelo
  \item Parámetros calibrados con datos reales chilenos
  \item Fuentes de datos y enlaces
  \item Análisis de sensibilidad e implicancias de política
\end{itemize}

% ============================================================
\section{Descripción del Sistema}
% ============================================================

\subsection{Límites del sistema}

El modelo considera la Región Metropolitana de Santiago como unidad
geográfica. El horizonte temporal relevante es de \textbf{1 a 10 años},
suficiente para capturar los efectos de demoras en infraestructura y
cambios en la flota.

\begin{definicion}
\textbf{Frontera del sistema:} Se incluyen los flujos de paquetes desde
centros de distribución hasta el consumidor final (última milla), la
flota de vehículos de reparto (motocicletas, furgones, vehículos eléctricos),
la infraestructura vial, los operadores logísticos y los consumidores.
Se excluyen: la cadena de abastecimiento primaria, el transporte
internacional y los servicios de mensajería no asociados a ecommerce.
\end{definicion}

\subsection{Variables principales}

La Tabla~\ref{tab:variables} resume las variables de estado (stocks),
tasas (flujos) y variables auxiliares del modelo.

\begin{table}[H]
\centering
\caption{Variables del modelo de última milla}
\label{tab:variables}
\begin{tabularx}{\textwidth}{@{}llXl@{}}
\toprule
\rowcolor{azulUCH!15}
\textbf{Variable} & \textbf{Tipo} & \textbf{Descripción} & \textbf{Unidad} \\
\midrule
$P(t)$ & Stock & Paquetes en tránsito en la RM & miles paquetes \\
$V(t)$ & Stock & Vehículos de reparto activos & vehículos \\
$C(t)$ & Stock & Congestión acumulada (índice) & adimensional \\
$D_e(t)$ & Stock & Demanda de ecommerce & millones USD/año \\
\midrule
$\lambda_P$ & Flujo entrada & Tasa de ingreso de paquetes & paquetes/día \\
$\mu_P$ & Flujo salida & Tasa de entrega exitosa & paquetes/día \\
$\phi_V$ & Flujo entrada & Incorporación de vehículos & veh./mes \\
$\delta_V$ & Flujo salida & Retiro/chatarrización de flota & veh./mes \\
\midrule
$\tau(C)$ & Auxiliar & Tiempo de entrega (función de $C$) & horas \\
$\kappa$ & Auxiliar & Costo por paquete & USD \\
$\rho$ & Auxiliar & Densidad de paquetes por vehículo & paq./veh. \\
$\alpha$ & Parámetro & Elasticidad demanda--precio & adimensional \\
\bottomrule
\end{tabularx}
\end{table}

% ============================================================
\section{Diagrama Causal (Bucles de Retroalimentación)}
% ============================================================

\subsection{Descripción de los bucles}

El sistema exhibe los siguientes bucles de retroalimentación
(\textbf{R} = refuerzo, \textbf{B} = balance):

\begin{description}[leftmargin=1cm,style=nextline]
  \item[\textcolor{rojoSD}{\textbf{R1 — Espiral de congestión}}]
    Mayor demanda de ecommerce $\rightarrow$ más paquetes $\rightarrow$
    más vehículos de reparto $\rightarrow$ mayor congestión $\rightarrow$
    mayores tiempos de entrega $\rightarrow$ operadores añaden más
    vehículos para cumplir SLA $\rightarrow$ mayor congestión
    (\textit{bucle de refuerzo positivo, amplificador}).

  \item[\textcolor{rojoSD}{\textbf{R2 — Efecto reputación / demanda}}]
    Mayores tiempos de entrega $\rightarrow$ peor experiencia del
    usuario $\rightarrow$ aumento de reseñas negativas $\rightarrow$
    mayor inversión en marketing para compensar $\rightarrow$
    mayor demanda de ecommerce (\textit{bucle de refuerzo débil}).

  \item[\textcolor{verdeSD}{\textbf{B1 — Precio como regulador}}]
    Mayor congestión $\rightarrow$ mayor costo de operación por
    paquete $\rightarrow$ alza del precio al consumidor $\rightarrow$
    reducción de la demanda de ecommerce (elasticidad negativa)
    $\rightarrow$ menos paquetes $\rightarrow$ menos vehículos
    (\textit{bucle de balance, estabilizador}).

  \item[\textcolor{verdeSD}{\textbf{B2 — Regulación municipal}}]
    Mayor congestión $\rightarrow$ presión regulatoria $\rightarrow$
    restricciones de acceso vehicular y zonas de carga $\rightarrow$
    reducción de vehículos en circulación $\rightarrow$
    menor congestión (\textit{bucle de balance con demora institucional
    de 1--3 años}).

  \item[\textcolor{verdeSD}{\textbf{B3 — Innovación modal}}]
    Alto costo de entrega $\rightarrow$ adopción de modos alternativos
    (bicicleta, casillero, drones) $\rightarrow$ reducción de
    vehículos motorizados $\rightarrow$ menor congestión.
\end{description}

\subsection{Diagrama causal (TikZ)}

La Figura~\ref{fig:causal} muestra el diagrama de relaciones causales.
Los arcos con \texttt{(+)} indican causalidad positiva (misma dirección)
y los arcos con \texttt{(--)} indican causalidad negativa (dirección opuesta).
Las etiquetas \textbf{R} y \textbf{B} identifican bucles.

\begin{figure}[H]
\centering
\begin{tikzpicture}[
  node distance = 2.2cm and 3.2cm,
  every node/.style = {align=center, font=\small},
  var/.style = {rectangle, rounded corners=4pt,
                minimum width=2.8cm, minimum height=0.9cm,
                draw=azulUCH!70, fill=azulUCH!8, thick},
  varR/.style = {rectangle, rounded corners=4pt,
                minimum width=2.8cm, minimum height=0.9cm,
                draw=rojoSD!70, fill=rojoSD!8, thick},
  varB/.style = {rectangle, rounded corners=4pt,
                minimum width=2.8cm, minimum height=0.9cm,
                draw=verdeSD!70, fill=verdeSD!8, thick},
  arr/.style = {-{Stealth[length=7pt]}, thick},
  arrN/.style = {-{Stealth[length=7pt]}, thick, dashed},
  lbl/.style = {font=\scriptsize\bfseries, inner sep=1pt},
  lblN/.style = {font=\scriptsize\bfseries, inner sep=1pt,
                 text=rojoSD},
]

% ── Nodos ────────────────────────────────────────────────────
\node[var]  (Dem)  at (0,0)      {Demanda\\ecommerce};
\node[var]  (Paq)  at (3.5,0)    {Paquetes\\en tránsito};
\node[varR] (Veh)  at (7,0)      {Vehículos de\\reparto activos};
\node[varR] (Con)  at (7,-3)     {Congestión\\urbana};
\node[var]  (Tie)  at (3.5,-3)   {Tiempo de\\entrega};
\node[varB] (Cos)  at (0,-3)     {Costo por\\paquete};
\node[varB] (Prec) at (0,-6)     {Precio al\\consumidor};
\node[varB] (Reg)  at (7,-6)     {Presión\\regulatoria};
\node[varB] (Alt)  at (3.5,-6)   {Modos\\alternativos};

% ── Arcos bucle R1 (sentido horario exterior) ─────────────────
\draw[arr, color=rojoSD!80]
  (Dem) -- node[above,lbl]{+} (Paq);
\draw[arr, color=rojoSD!80]
  (Paq) -- node[above,lbl]{+} (Veh);
\draw[arr, color=rojoSD!80]
  (Veh) -- node[right,lbl]{+} (Con);
\draw[arr, color=rojoSD!80]
  (Con) -- node[above,lbl]{+} (Tie);
\draw[arr, color=rojoSD!80]
  (Tie) to[bend left=25] node[left,lbl]{+} (Veh);

% ── Arcos bucle B1 (precio) ───────────────────────────────────
\draw[arr, color=verdeSD!80]
  (Con) -- node[left,lbl]{+} (Cos);
\draw[arr, color=verdeSD!80]
  (Cos) -- node[left,lbl]{+} (Prec);
\draw[arrN, color=verdeSD!80]
  (Prec) -- node[below,lbl]{--} (Dem);

% ── Arcos bucle B2 (regulación) ───────────────────────────────
\draw[arr, color=ambarSD!80]
  (Con) -- node[right,lbl]{+} (Reg);
\draw[arrN, color=ambarSD!80]
  (Reg) to[bend right=20] node[right,lbl,xshift=3pt]{--} (Veh);

% ── Arcos bucle B3 (modos alternativos) ──────────────────────
\draw[arr, color=verdeSD!60]
  (Cos) -- node[below,lbl]{+} (Alt);
\draw[arrN, color=verdeSD!60]
  (Alt) -- node[below,lbl]{--} (Con);

% ── Tiempo de entrega → demanda (R2 débil) ───────────────────
\draw[arrN, color=rojoSD!40, dashed]
  (Tie) to[bend right=30] node[below,lbl,text=rojoSD!60]{--} (Dem);

% ── Etiquetas de bucle ────────────────────────────────────────
\node[font=\bfseries\small,text=rojoSD] at (5.3,-1.5) {R1};
\node[font=\bfseries\small,text=rojoSD] at (1.7,-1.5) {R2};
\node[font=\bfseries\small,text=verdeSD] at (-1.2,-3) {B1};
\node[font=\bfseries\small,text=ambarSD] at (8.5,-3) {B2};
\node[font=\bfseries\small,text=verdeSD] at (1.7,-5.5) {B3};

\end{tikzpicture}
\caption{Diagrama causal del sistema de última milla en Santiago.
  Arcos sólidos~(+): causalidad positiva.
  Arcos discontinuos~(--): causalidad negativa.
  \textcolor{rojoSD}{\textbf{Rojo}}: bucles de refuerzo;
  \textcolor{verdeSD}{\textbf{verde}}: bucles de balance;
  \textcolor{ambarSD}{\textbf{ámbar}}: regulación.}
\label{fig:causal}
\end{figure}

% ============================================================
\section{Diagrama de Stocks y Flujos}
% ============================================================

\subsection{Estructura}

El modelo tiene \textbf{4 stocks principales} conectados por flujos
controlados por variables auxiliares.

\begin{figure}[H]
\centering
\begin{tikzpicture}[
  stock/.style = {rectangle, minimum width=2.6cm, minimum height=1.1cm,
                  draw=azulUCH, fill=azulUCH!10, thick, align=center, font=\small\bfseries},
  flow/.style  = {-{Stealth[length=9pt]}, ultra thick, color=azulUCH!70},
  cloud/.style = {ellipse, minimum width=1.5cm, minimum height=0.8cm,
                  draw=gray!60, fill=gray!8, font=\scriptsize, align=center},
  aux/.style   = {circle, minimum size=0.7cm,
                  draw=verdeSD!70, fill=verdeSD!8, font=\scriptsize, align=center},
  info/.style  = {-{Stealth[length=5pt]}, dashed, color=grisOscuro, thin},
  valve/.style = {regular polygon, regular polygon sides=3,
                  draw=azulUCH, fill=white, minimum size=0.5cm},
  every node/.style={align=center},
]

% ── Row 1: Demanda ───────────────────────────────────────────
\node[cloud]  (SrcD) at (-3.5, 0) {};
\node[stock]  (Dem)  at (0,    0) {$D_e(t)$\\Demanda\\ecommerce};
\node[cloud]  (SnkD) at (3.5,  0) {};

\draw[flow] (SrcD) -- node[above,font=\scriptsize]{$g_D$} (Dem);
\draw[flow] (Dem)  -- node[above,font=\scriptsize]{$d_D$} (SnkD);

% ── Row 2: Paquetes ───────────────────────────────────────────
\node[cloud]  (SrcP) at (-3.5,-2.5) {};
\node[stock]  (Paq)  at (0,   -2.5) {$P(t)$\\Paquetes\\en tránsito};
\node[cloud]  (SnkP) at (3.5, -2.5) {};

\draw[flow] (SrcP) -- node[above,font=\scriptsize]{$\lambda_P$} (Paq);
\draw[flow] (Paq)  -- node[above,font=\scriptsize]{$\mu_P$}     (SnkP);

% ── Row 3: Vehículos ──────────────────────────────────────────
\node[cloud]  (SrcV) at (-3.5,-5) {};
\node[stock]  (Veh)  at (0,   -5) {$V(t)$\\Vehículos\\activos};
\node[cloud]  (SnkV) at (3.5, -5) {};

\draw[flow] (SrcV) -- node[above,font=\scriptsize]{$\phi_V$} (Veh);
\draw[flow] (Veh)  -- node[above,font=\scriptsize]{$\delta_V$} (SnkV);

% ── Row 4: Congestión ─────────────────────────────────────────
\node[cloud]  (SrcC) at (-3.5,-7.5) {};
\node[stock]  (Con)  at (0,   -7.5) {$C(t)$\\Congestión\\(índice)};
\node[cloud]  (SnkC) at (3.5, -7.5) {};

\draw[flow] (SrcC) -- node[above,font=\scriptsize]{$\beta_C$} (Con);
\draw[flow] (Con)  -- node[above,font=\scriptsize]{$\gamma_C$} (SnkC);

% ── Variables auxiliares ──────────────────────────────────────
\node[aux] (Tau) at (5.5,-3.8) {$\tau(C)$\\tiempo};
\node[aux] (Kap) at (5.5,-6)   {$\kappa$\\costo};
\node[aux] (Rho) at (-5.5,-3.8){$\rho$\\densidad};

% ── Vínculos informativos ─────────────────────────────────────
\draw[info] (Dem)  -- node[right,font=\tiny]{conv.} (Paq);
\draw[info] (Paq)  -- node[right,font=\tiny]{necesidad} (Veh);
\draw[info] (Veh)  -- node[right,font=\tiny]{genera} (Con);
\draw[info] (Con)  to[bend right=20] (Tau);
\draw[info] (Tau)  to[bend right=20] (Veh);
\draw[info] (Con)  to[bend left=20]  (Kap);
\draw[info] (Kap)  to[bend left=30]  (Dem);
\draw[info] (Paq)  to[bend left=10]  (Rho);
\draw[info] (Rho)  to[bend left=10]  (Veh);

\end{tikzpicture}
\caption{Diagrama de stocks y flujos. Las nubes ($\sim$) representan
  fuentes/sumideros fuera del límite del sistema.
  Las flechas discontinuas son vínculos de información (no flujos de materia).}
\label{fig:stocks}
\end{figure}

% ============================================================
\section{Arquetipos Sistémicos}
% ============================================================

\begin{arquetipobox}{Límites al Crecimiento}
\textbf{Descripción:} El ecommerce crece impulsado por la conveniencia
(bucle R1), pero la capacidad de la red vial actúa como restricción.
A medida que la congestión aumenta, el tiempo de entrega crece y la
experiencia del usuario se deteriora, frenando el crecimiento de la demanda
(bucle B1).

\textbf{Síntoma en datos:} La tasa de crecimiento del ecommerce chileno
bajó del 30\,\% anual (2019--2021) al 12\,\% (2022--2023), coincidiendo
con el aumento de tiempos de entrega reportados por operadores.

\textbf{Palanca sistémica:} No aumentar la flota, sino redistribuir la
demanda en el tiempo (ventanas de entrega programadas) y en el espacio
(microbodegas en comunas periféricas).
\end{arquetipobox}

\vspace{0.5em}

\begin{arquetipobox}{Soluciones que Fallan (\textit{Fixes that Fail})}
\textbf{Descripción:} Ante las demoras, los operadores logísticos añaden
vehículos como solución de corto plazo. Esto alivia temporalmente el
tiempo de entrega, pero a mediano plazo intensifica la congestión,
empeorando el problema original.

\textbf{Síntoma en datos:} En Santiago, la cantidad de vehículos de
reparto creció un 65\,\% entre 2020 y 2023 (SECTRA), mientras el índice
de congestión TomTom para Santiago subió de 35\,\% (2020) a 47\,\% (2023).

\textbf{Palanca sistémica:} Invertir en modelos de consolidación de carga
(City Logistics Hubs) en lugar de aumentar la flota individual.
\end{arquetipobox}

\vspace{0.5em}

\begin{arquetipobox}{Tragedia de los Comunes (\textit{Tragedy of the Commons})}
\textbf{Descripción:} Cada operador logístico individualmente maximiza
su uso de la red vial, que es un recurso compartido. La suma de decisiones
racionales individuales lleva a una congestión que perjudica a todos,
incluyendo a los propios operadores.

\textbf{Síntoma en datos:} Santiago tiene más de 15 operadores logísticos
de ecommerce relevantes (Correos de Chile, Starken, Chilexpress, BlueExpress,
Shippify, etc.) compitiendo por la misma infraestructura vial.

\textbf{Palanca sistémica:} Coordinación de rutas y ventanas horarias
mediante plataforma regulada (similar al modelo de París o Ámsterdam).
\end{arquetipobox}

\vspace{0.5em}

\begin{arquetipobox}{Escalada (\textit{Escalation})}
\textbf{Descripción:} Los operadores compiten por velocidad de entrega
(same-day, same-hour), lo que incentiva la adición de más vehículos y
más viajes por paquete, escalando la congestión mutuamente.

\textbf{Síntoma en datos:} MercadoLibre, Falabella y Amazon Chile
compiten en entrega en el mismo día, lo que requiere vehículos dedicados
en horarios pico.

\textbf{Palanca sistémica:} Estándar sectorial de tiempos de entrega
mínimos razonables o tarificación por congestión.
\end{arquetipobox}

% ============================================================
\section{Ecuaciones Diferenciales del Modelo}
% ============================================================

\subsection{Sistema de EDOs}

El modelo se formula como un sistema de cuatro ecuaciones diferenciales
ordinarias de primer orden, más ecuaciones auxiliares.

\subsubsection{Stock 1: Demanda de ecommerce $D_e(t)$}

\begin{equation}
\boxed{
\frac{dD_e}{dt} = r_D \cdot D_e(t)\left(1 - \frac{D_e(t)}{K_D}\right)
  - \alpha \cdot \kappa(t) \cdot D_e(t)
  - \beta_{D} \cdot \tau(t)^{\theta}
}
\label{eq:demanda}
\end{equation}

donde:
\begin{itemize}[leftmargin=2cm,itemsep=2pt]
  \item $r_D \approx 0.25\ \text{año}^{-1}$: tasa de crecimiento orgánico del ecommerce
  \item $K_D$: capacidad de saturación del mercado (estimada en 3 veces el nivel actual)
  \item $\alpha \approx -0.8$: elasticidad precio--demanda de ecommerce en Chile (CCS, 2023)
  \item $\kappa(t)$: costo de envío normalizado (ver ecuación~\ref{eq:costo})
  \item $\beta_D \approx 0.05$: sensibilidad de la demanda al tiempo de entrega
  \item $\theta \approx 1.2$: exponente de aversión a las demoras
\end{itemize}

\subsubsection{Stock 2: Paquetes en tránsito $P(t)$}

\begin{equation}
\boxed{
\frac{dP}{dt} = \lambda_P(t) - \mu_P(t)
}
\label{eq:paquetes}
\end{equation}

\begin{equation}
\lambda_P(t) = \eta \cdot D_e(t) \qquad\qquad
\mu_P(t) = \frac{P(t)}{\tau(t)}
\end{equation}

donde:
\begin{itemize}[leftmargin=2cm,itemsep=2pt]
  \item $\lambda_P$: tasa de ingreso de paquetes al sistema ($\eta \approx 1{,}8\ \text{paq./USD}$)
  \item $\mu_P$: tasa de entrega, inversamente proporcional al tiempo de entrega $\tau(t)$
  \item $\tau(t)$: tiempo promedio de entrega (horas), función del nivel de congestión
\end{itemize}

\subsubsection{Stock 3: Vehículos de reparto activos $V(t)$}

\begin{equation}
\boxed{
\frac{dV}{dt} = \phi_V(t) - \delta_V \cdot V(t)
}
\label{eq:vehiculos}
\end{equation}

La tasa de incorporación de vehículos depende del desequilibrio entre
la demanda de cobertura y la capacidad actual de la flota:

\begin{equation}
\phi_V(t) = \frac{1}{T_V}\max\!\left(0,\;
  \underbrace{\frac{P(t)}{\rho_0}}_{\text{flota requerida}}
  - V(t)\right)
\end{equation}

donde:
\begin{itemize}[leftmargin=2cm,itemsep=2pt]
  \item $T_V \approx 3\ \text{meses}$: demora de ajuste de la flota
    (tiempo para contratar, habilitar y desplegar un vehículo/conductor)
  \item $\rho_0 \approx 40\ \text{paq./veh./día}$: productividad estándar
    de un vehículo de reparto en Santiago (TomTom + estimación operadores)
  \item $\delta_V \approx 0.15\ \text{año}^{-1}$: tasa de retiro/renovación de flota
\end{itemize}

\subsubsection{Stock 4: Índice de congestión $C(t)$}

\begin{equation}
\boxed{
\frac{dC}{dt} = \beta_C \cdot f\!\left(\frac{V(t)}{V_{\max}}\right) - \gamma_C \cdot C(t)
  - \xi \cdot R(t)
}
\label{eq:congestion}
\end{equation}

La función de generación de congestión usa la fórmula del Bureau of
Public Roads (BPR), estándar en ingeniería de transporte:

\begin{equation}
f\!\left(\frac{V}{V_{\max}}\right) = 1 + 0.15\left(\frac{V(t)}{V_{\max}}\right)^4
\end{equation}

donde:
\begin{itemize}[leftmargin=2cm,itemsep=2pt]
  \item $V_{\max}$: capacidad vial máxima de la RM en vehículos equivalentes
    ($\approx 2{,}1 \times 10^6\ \text{veh.}$ según SECTRA 2023)
  \item $\gamma_C \approx 0.3\ \text{día}^{-1}$: tasa de disipación de congestión
    (relacionada con la longitud promedio de viaje y velocidad libre)
  \item $\xi$: efecto de la regulación $R(t)$ sobre la congestión
\end{itemize}

\subsection{Ecuaciones auxiliares}

\subsubsection{Tiempo de entrega $\tau(t)$}

\begin{equation}
\tau(t) = \tau_0 \cdot \left[1 + 0.15\left(\frac{V(t)}{V_{\max}}\right)^4\right]
\label{eq:tau}
\end{equation}

con $\tau_0 \approx 4.5\ \text{h}$ (tiempo de entrega en condiciones sin congestión,
estimado a partir de datos de operadores en Santiago).

\subsubsection{Costo por paquete $\kappa(t)$}

\begin{equation}
\kappa(t) = \kappa_0 + c_V \cdot V(t) \cdot \tau(t)
\label{eq:costo}
\end{equation}

donde $\kappa_0 \approx 2.500\ \text{CLP}$ es el costo base y $c_V$ es el costo
variable por hora--vehículo ($\approx 800\ \text{CLP/h}$, incluye combustible,
depreciación y remuneración mínima del conductor).

\subsubsection{Regulación $R(t)$ (dinámica con demora)}

\begin{equation}
\frac{dR}{dt} = \frac{1}{T_R}\left[R^*(C(t)) - R(t)\right]
\label{eq:regulacion}
\end{equation}

con $T_R \approx 18\ \text{meses}$ (demora institucional para implementar regulación)
y $R^*(C) = R_{\max}\cdot\dfrac{C(t)}{C_{1/2} + C(t)}$ (función de Hill/saturación).

\subsection{Resumen del sistema}

El sistema completo es:

\begin{equation}
\begin{cases}
\dfrac{dD_e}{dt} = r_D D_e\!\left(1 - \dfrac{D_e}{K_D}\right)
  - \alpha\,\kappa\, D_e - \beta_D\,\tau^\theta \\[12pt]
\dfrac{dP}{dt}   = \eta\,D_e - \dfrac{P}{\tau} \\[10pt]
\dfrac{dV}{dt}   = \dfrac{1}{T_V}\max\!\left(0,\dfrac{P}{\rho_0}-V\right) - \delta_V V \\[10pt]
\dfrac{dC}{dt}   = \beta_C\!\left[1+0.15\!\left(\tfrac{V}{V_{\max}}\right)^4\right]
  - \gamma_C C - \xi\,R \\[10pt]
\dfrac{dR}{dt}   = \dfrac{1}{T_R}\!\left[R^*(C) - R\right]
\end{cases}
\label{eq:sistema}
\end{equation}

\subsection{Análisis de puntos de equilibrio}

En el punto de equilibrio ($\dot{D}_e = \dot{P} = \dot{V} = \dot{C} = \dot{R} = 0$):

\begin{align}
P^* &= \eta\,D_e^*\,\tau^* \\
V^* &= \frac{P^*}{\rho_0} = \frac{\eta\,D_e^*\,\tau^*}{\rho_0} \\
C^* &= \frac{\beta_C}{\gamma_C}\!\left[1+0.15\!\left(\frac{V^*}{V_{\max}}\right)^4\right]
  - \frac{\xi\,R^*}{\gamma_C}
\end{align}

La estabilidad del equilibrio se analiza con la matriz Jacobiana
$J = \partial f_i / \partial x_j$ evaluada en $(D_e^*, P^*, V^*, C^*, R^*)$.
El sistema exhibe equilibrios múltiples: un equilibrio de ``baja congestión''
estable y un equilibrio de ``alta congestión'' también localmente estable,
separados por un umbral inestable — característico del arquetipo
\textit{Límites al Crecimiento} con efectos de umbral.

% ============================================================
\section{Parámetros del Modelo: Datos Reales Chile}
% ============================================================

\begin{databox}
Los parámetros a continuación fueron estimados a partir de fuentes
oficiales chilenas y publicaciones sectoriales. Se incluyen los URLs
de acceso directo.
\end{databox}

\subsection{Tabla de parámetros calibrados}

\begin{longtable}{@{}p{2cm}p{2.5cm}p{2.5cm}p{1.5cm}p{4cm}@{}}
\caption{Parámetros del modelo y fuentes de calibración}
\label{tab:parametros}\\
\toprule
\rowcolor{azulUCH!15}
\textbf{Param.} & \textbf{Valor} & \textbf{Rango} & \textbf{Unidad}
  & \textbf{Fuente} \\
\midrule
\endfirsthead
\multicolumn{5}{c}{\small\textit{(continuación de la tabla)}}\\
\toprule
\rowcolor{azulUCH!15}
\textbf{Param.} & \textbf{Valor} & \textbf{Rango} & \textbf{Unidad}
  & \textbf{Fuente} \\
\midrule
\endhead
\midrule
\multicolumn{5}{r}{\small\textit{Continúa en la página siguiente}}\\
\endfoot
\bottomrule
\endlastfoot

$r_D$       & 0.25   & 0.10--0.35  & año$^{-1}$
  & CCS, Reporte eCommerce Chile 2023 \\

$\alpha$    & --0.8  & --0.5 a --1.2 & adim.
  & Banco Central, Informe IPC y consumo privado \\

$\tau_0$    & 4.5    & 3--8        & horas
  & TomTom Traffic Index Santiago 2023 \\

$\rho_0$    & 40     & 25--55      & paq./veh./día
  & Estimación operadores (Starken, Chilexpress anuales) \\

$T_V$       & 3      & 2--6        & meses
  & Entrevistas sectoriales / SECTRA \\

$\delta_V$  & 0.15   & 0.10--0.20  & año$^{-1}$
  & SII (Servicio de Impuestos Internos), padrón vehicular \\

$V_{\max}$  & $2.1\times10^6$ & -- & veh. equiv.
  & SECTRA, Encuesta Origen--Destino RM 2021 \\

$\beta_C$   & 0.012  & 0.008--0.02 & día$^{-1}$
  & TomTom Traffic Index + SECTRA \\

$\gamma_C$  & 0.30   & 0.20--0.45  & día$^{-1}$
  & Calibración TomTom hora pico vs. hora valle \\

$T_R$       & 18     & 12--36      & meses
  & Análisis de ordenanzas municipales Santiago \\

$\kappa_0$  & 2\,500 & 2\,000--4\,000 & CLP
  & SERNAC, análisis de tarifas publicadas (2023) \\

$c_V$       & 800    & 600--1\,200 & CLP/h
  & Dirección del Trabajo, remuneraciones sector transporte \\

$\eta$      & 1.8    & 1.4--2.5    & paq./USD
  & CCS + Correos de Chile informe anual \\

\end{longtable}

\subsection{Condiciones iniciales (2023)}

\begin{align}
D_e(0) &= 8{,}900 \text{ millones USD/año}
  \quad \text{(CCS, Reporte eCommerce 2023)} \\
P(0)   &= 350{,}000 \text{ paquetes en tránsito}
  \quad \text{(estimado, Correos + operadores privados)} \\
V(0)   &= 28{,}000 \text{ vehículos de reparto}
  \quad \text{(SECTRA 2023, vehículos livianos comerciales)} \\
C(0)   &= 0.47 \text{ (índice TomTom = 47\,\% sobre tiempo libre)}
  \quad \text{(TomTom 2023)} \\
R(0)   &= 0.20 \text{ (nivel regulatorio actual, escala 0--1)}
\end{align}

% ============================================================
\section{Fuentes de Datos y Acceso}
% ============================================================

\begin{databox}
Todas las fuentes son de acceso gratuito. Los URLs fueron verificados
en mayo de 2025.
\end{databox}

\subsection{Fuentes nacionales}

\begin{enumerate}[leftmargin=1.5cm,label=\textbf{[\arabic*]}]

\item \textbf{SECTRA — Subsecretaría de Transportes, Encuesta Origen--Destino
  Región Metropolitana 2021}\\
  Datos de flujos vehiculares por hora, comuna, tipo de vehículo.\\
  \url{https://www.sectra.gob.cl/encuestas_movilidad/encuestas_movilidad.htm}

\item \textbf{Instituto Nacional de Estadísticas (INE) — Encuesta de
  Comercio y Servicios (ECOM)}\\
  Ventas del comercio minorista con desglose por canal (presencial/online).\\
  \url{https://www.ine.gob.cl/estadisticas/economia/comercio-servicios-y-turismo/comercio-interno}

\item \textbf{Cámara de Comercio de Santiago (CCS) — Reporte Anual
  eCommerce Chile}\\
  Facturación, crecimiento, categorías, medios de pago y logística.\\
  \url{https://www.ccs.cl/wp-content/uploads/2023/06/Reporte-eCommerce-Chile-2023.pdf}

\item \textbf{SERNAC — Estadísticas de Reclamos por Sector}\\
  Reclamos de consumidores sobre tiempos de entrega y comercio electrónico.\\
  \url{https://www.sernac.cl/portal/604/w3-propertyvalue-29834.html}

\item \textbf{Banco Central de Chile — Base de Datos Estadísticos (BDE)}\\
  IPC, consumo privado, tasa de cambio. Datos descargables en Excel/API.\\
  \url{https://si3.bcentral.cl/siete}

\item \textbf{Dirección del Trabajo — Encuesta Laboral (ENCLA)}\\
  Remuneraciones por sector, incluyendo transporte y logística.\\
  \url{https://www.dt.gob.cl/portal/1629/w3-propertyvalue-22936.html}

\item \textbf{Ministerio de Hacienda — Portal Datos Abiertos (datos.gob.cl)}\\
  Repositorio central de datasets gubernamentales.\\
  \url{https://datos.gob.cl}

\item \textbf{Servicio de Impuestos Internos (SII) — Estadísticas de
  Empresas por Rubro}\\
  Empresas activas en transporte de carga y mensajería (CIIU 5320).\\
  \url{https://www.sii.cl/estadisticas/empresas.htm}

\end{enumerate}

\subsection{Fuentes internacionales con datos para Santiago}

\begin{enumerate}[leftmargin=1.5cm,label=\textbf{[\arabic*]},resume]

\item \textbf{TomTom Traffic Index — Santiago de Chile}\\
  Índice de congestión histórico, comparación hora pico/valle, tendencia anual.\\
  \url{https://www.tomtom.com/traffic-index/santiago-traffic/}

\item \textbf{ITF/OECD — Urban Freight and Logistics}\\
  Benchmarks internacionales de última milla, indicadores de eficiencia.\\
  \url{https://www.itf-oecd.org/urban-freight-logistics}

\item \textbf{World Bank — Logistics Performance Index (LPI)}\\
  Chile ocupa el puesto 34 mundial (2023), con subíndices de puntualidad.\\
  \url{https://lpi.worldbank.org/international/scorecard/line/386/C/CHL}

\item \textbf{Correos de Chile — Memoria Anual}\\
  Volúmenes de envíos, tiempos de entrega declarados, red de distribución.\\
  \url{https://www.correos.cl/web/guest/nuestra-empresa/memorias-anuales}

\end{enumerate}

% ============================================================
\section{Análisis de Sensibilidad}
% ============================================================

Los parámetros críticos — aquellos cuya variación más afecta el
estado de equilibrio — son, en orden de importancia:

\begin{enumerate}[leftmargin=1.5cm]
  \item \textbf{$T_V$ (demora de ajuste de flota)}: Una reducción de
    3 a 1 mes (por ejemplo, mediante plataformas de gig economy)
    amplifica el bucle R1 y lleva al sistema más rápido al equilibrio
    de alta congestión.

  \item \textbf{$\rho_0$ (productividad por vehículo)}: Aumentar
    $\rho_0$ de 40 a 70 paq./día mediante optimización de rutas
    (IA, consolidación) reduce $V^*$ en un 43\,\%, el efecto de
    política más potente disponible.

  \item \textbf{$T_R$ (demora regulatoria)}: Reducir la demora de
    implementación de regulación de 18 a 6 meses acelera la activación
    del bucle B2, reduciendo el pico de congestión en un 18\,\% según
    simulaciones preliminares.

  \item \textbf{$\alpha$ (elasticidad precio--demanda)}: La elasticidad
    chilena de --0.8 es menor (en valor absoluto) que la europea
    ($\approx -1.2$), lo que implica que el mecanismo de precio (B1)
    es menos efectivo como estabilizador en el contexto local.
\end{enumerate}

% ============================================================
\section{Implicancias de Política}
% ============================================================

El análisis sistémico sugiere tres palancas de intervención,
ordenadas por efectividad relativa:

\begin{table}[H]
\centering
\caption{Palancas de política y bucle que activan}
\label{tab:politica}
\begin{tabularx}{\textwidth}{@{}p{3.5cm}Xp{2.5cm}p{2cm}@{}}
\toprule
\rowcolor{azulUCH!15}
\textbf{Intervención} & \textbf{Mecanismo} & \textbf{Bucle} & \textbf{Horizonte} \\
\midrule
Hubs de consolidación urbana
  & Aumenta $\rho_0$, reduce $V^*$ sin restringir demanda
  & B1, B3 & 2--4 años \\
\midrule
Regulación de ventanas horarias
  & Redistribuye la carga, reduce el pico de $C(t)$
  & B2 & 1--2 años \\
\midrule
Tarificación por congestión (congestion pricing)
  & Introduce costo variable que refuerza B1
  & B1 & 3--5 años \\
\midrule
Subsidio a bicicletas de carga y vehículos eléctricos
  & Reduce $c_V$, aumenta $\rho_0$, activa B3
  & B3 & 1--3 años \\
\bottomrule
\end{tabularx}
\end{table}

\textbf{Trampa sistémica a evitar:} Intervenciones que solo aumenten
la capacidad vial (más pistas) activan el efecto rebote
(\textit{induced demand}), generando tráfico latente que llena
la nueva capacidad y devuelve el sistema al equilibrio de alta
congestión — este es el arquetipo \textit{Soluciones que Fallan}
aplicado a infraestructura.

% ============================================================
\section{Conclusiones}
% ============================================================

El sistema de última milla en Santiago presenta una estructura dinámica
dominada por un bucle de refuerzo (R1) que tiende a amplificar la
congestión, contrarrestado por tres bucles de balance de distinta
velocidad: el precio (B1, rápido), la regulación (B2, lento, 18 meses
de demora) y la innovación modal (B3, mediano plazo).

Los cuatro arquetipos identificados — Límites al Crecimiento, Soluciones
que Fallan, Tragedia de los Comunes y Escalada — señalan que las
intervenciones de corto plazo (más vehículos, más pistas) son
contraproducentes. Las palancas de alto apalancamiento se encuentran
en el aumento de la productividad por vehículo ($\rho_0$) y en la
reducción de la demora regulatoria ($T_R$).

El modelo EDO completo (ecuación~\ref{eq:sistema}) puede ser simulado
en Vensim, Stella Architect, Python (\texttt{pysd}) o R con el paquete
\texttt{deSolve}, usando los parámetros y condiciones iniciales
reportados en las Secciones 6 y 7.

% ============================================================
\appendix
\section{Código de Simulación (Python / pysd)}
% ============================================================

\begin{lstlisting}[language=Python, caption={Simulación básica del
  sistema con \texttt{scipy.integrate.solve\_ivp}}]
import numpy as np
from scipy.integrate import solve_ivp
import matplotlib.pyplot as plt

# Parametros
r_D    = 0.25;   K_D    = 26700    # millones USD
alpha  = -0.8;   beta_D = 0.05
eta    = 1.8;    tau0   = 4.5
rho0   = 40;     T_V    = 0.25     # años
delta_V= 0.15;   V_max  = 2.1e6
beta_C = 0.012 * 365  # convertir a año^-1
gamma_C= 0.30 * 365;  xi = 0.1
T_R    = 1.5;    R_max  = 0.8;     C_half = 0.5
kappa0 = 2500;   c_V    = 800 * 24 * 365  # anual
theta  = 1.2

def tau(V):
    return tau0 * (1 + 0.15 * (V / V_max)**4)

def kappa(V):
    t = tau(V)
    return kappa0 + c_V * V * t

def R_star(C):
    return R_max * C / (C_half + C)

def sistema(t, y):
    De, P, V, C, R = y
    t_val = tau(V)
    k_val = kappa(V)

    dDe = r_D*De*(1 - De/K_D) - alpha*k_val*De - beta_D*t_val**theta
    dP  = eta*De - P/t_val
    dV  = max(0, P/rho0 - V) / T_V - delta_V*V
    dC  = beta_C*(1 + 0.15*(V/V_max)**4) - gamma_C*C - xi*R
    dR  = (R_star(C) - R) / T_R
    return [dDe, dP, dV, dC, dR]

# Condiciones iniciales 2023
y0 = [8900, 350000, 28000, 0.47, 0.20]
t_span = (0, 10)
t_eval = np.linspace(0, 10, 500)

sol = solve_ivp(sistema, t_span, y0, t_eval=t_eval, method='RK45')

fig, axes = plt.subplots(2, 2, figsize=(12, 8))
labels = ['Demanda ecommerce (M USD)', 'Paquetes en transito',
          'Vehiculos activos', 'Indice de congestion']
for i, ax in enumerate(axes.flatten()):
    ax.plot(sol.t, sol.y[i])
    ax.set_xlabel('Años desde 2023')
    ax.set_title(labels[i])
    ax.grid(True, alpha=0.3)
plt.tight_layout()
plt.savefig('simulacion_ultima_milla.png', dpi=150)
plt.show()
\end{lstlisting}

% ============================================================
\section{Glosario}
% ============================================================

\begin{description}[leftmargin=3cm,style=nextline,font=\bfseries]
  \item[Bucle de refuerzo (R)]
    Ciclo de retroalimentación donde un aumento en una variable
    genera, a través de la cadena causal, un aumento adicional
    en la misma variable (retroalimentación positiva, amplificadora).
  \item[Bucle de balance (B)]
    Ciclo de retroalimentación donde un aumento en una variable
    genera, a través de la cadena causal, una disminución posterior
    en la misma variable (retroalimentación negativa, estabilizadora).
  \item[Stock]
    Acumulación en el sistema que cambia a través de flujos.
    Corresponde a la variable de estado en las EDOs.
  \item[Flujo]
    Tasa de cambio de un stock (derivada temporal).
  \item[Arquetipo sistémico]
    Patrón de estructura de bucles que produce un comportamiento
    dinámico reconocible, descrito originalmente por Senge (1990).
  \item[Última milla]
    Segmento final de la cadena logística desde el centro de
    distribución hasta el punto de entrega al cliente final.
  \item[TEU]
    Twenty-foot Equivalent Unit, unidad estándar de medida
    de contenedores marítimos.
  \item[TomTom Traffic Index]
    Índice internacional que mide el exceso de tiempo de viaje
    respecto a condiciones de tráfico libre, en porcentaje.
\end{description}

\end{document}
